% Generated draft: requires mathematical and source review.
\documentclass[11pt]{article}
\usepackage[margin=25mm]{geometry}
\usepackage{fontspec}
\setmainfont{Latin Modern Roman}
\usepackage{amsmath,amssymb,mathtools,unicode-math}
\setmathfont{Latin Modern Math}
\usepackage{enumitem,needspace,xurl,hyperref}
\setlength{\parindent}{0pt}
\setlength{\parskip}{5pt}
\setlength{\emergencystretch}{3em}
\newcommand{\N}{\mathbb{N}}
\newcommand{\Z}{\mathbb{Z}}
\newcommand{\Q}{\mathbb{Q}}
\newcommand{\R}{\mathbb{R}}
\newcommand{\C}{\mathbb{C}}
\newcommand{\F}{\mathbb{F}}
\newcommand{\A}{\mathbb{A}}
\newcommand{\PP}{\mathbb{P}}
\newcommand{\E}{\mathbb{E}}
\newcommand{\eps}{\varepsilon}
\newcommand{\dd}{\,\mathrm{d}}
\begin{document}
\section*{UnsolvedMath ID: 2309006}
\textbf{Source problem code:} catalogue entry
\subsection{Definitions and mathematical statement}
Let $\widehat\C=\C\cup\{\infty\}$, let $D\subset\widehat\C$ be open with $\partial D\ne\varnothing$, and let $F\subset D$ be relatively closed. Suppose $f:F\to\C$ is uniformly approximated by functions $f_n$ analytic on $D$:
\[
 \sup_{z\in F}|f_n(z)-f(z)|\longrightarrow0.
\]
The source then assumes that $f$ extends continuously to $F\cup E$ and asks whether each $f_n$ can also be extended continuously there.

The printed location of $E$ is
\[
 E\subseteq\partial F_n\cap\partial D,
\]
but no sets $F_n$ are introduced. If the intended symbol is $\partial F$, that must be the boundary in the ambient sphere, not merely the relative boundary inside $D$, since the latter cannot meet $\partial D$. This is a possible correction, not one silently imposed.

There is also an essential quantifier issue. The literal reading asks for continuous extensions of the already given restrictions $f_n|_F$. A different approximation problem asks whether the approximating sequence can be chosen or replaced by analytic functions whose restrictions extend continuously to $F\cup E$. Uniform convergence to boundary-continuous data does not by itself require each approximating term to have that boundary continuity. The source's asserted positive special cases may concern such a choice/replacement theorem, but no wording specifying it is supplied. Requiring an extension on $D\cup E$, rather than only of the restriction to $F\cup E$, is another distinct assertion.

The source reports an affirmative answer when $D$ is the unit disk or $E$ is compact. These are retained as reported source claims, with the unresolved set notation, extension domain and quantifiers explicit; they are not promoted to a theorem under the literal ambiguous reading.
\subsection{Short English statement}
Ask whether analytic approximation on a relatively closed set can preserve continuous boundary extension at the specified boundary points. The source contains an undefined $F_n$ and does not distinguish extension of an arbitrary given approximant from the ability to select suitable approximants. Those distinctions are necessary to reconcile the reported disk and compact-$E$ cases. The supplied triage reports the problem open in Hayman's 2018 update and no definitive resolution located; no independent status review is asserted.
\subsection{Sources}
\begin{itemize}
\item UnsolvedMath permanent ID 2309006, source code catalogue entry, ``Research Problems in Function Theory --- Problem 9.6.'' Registry curation and metadata: Ulam AI and UnsolvedMath Contributors, CC BY 4.0; updated 2026-08-06.
\item Hayman, \emph{Research Problems in Function Theory} (2018 source update), Problem 9.6, attributed to A. Stray: \url{https://arxiv.org/abs/1809.07200}. Supplied source-TeX extraction accessed 2026-07-29. The metadata retains the update for status triage and marks status/release and rights review NEEDS\_REVIEW.
\item No separate complete reference is supplied for the unit-disk or compact-boundary-set positive cases. No missing theorem or correction is inferred.
\end{itemize}
\end{document}
